% The template every paper starts from. Packages used here all ship with the pinned
% TeX Live 2022 (TinyTeX); adding one is documented in docs/typesetting.md.
\documentclass[11pt]{article}
\usepackage[margin=1in]{geometry}
\usepackage{amsmath,graphicx,booktabs,xcolor,url}
\usepackage[numbers]{natbib}
\usepackage[hidelinks]{hyperref}
\graphicspath{{build/figures/}}

\title{Template: a paper that builds from sources}
\author{Authors decided by the user}
\date{Draft, not for distribution}

\begin{document}
\maketitle

\begin{abstract}
This template proves the build: the PDF is made from \texttt{main.tex} and \texttt{refs.bib},
and every figure is regenerated from its script, never committed.
\end{abstract}

\section{Introduction}
Every claim traces to evidence: a commit, a test, a reproduced measurement, or a cited
source~\citep{knuth1984texbook}. Figure~\ref{fig:hello} is generated by
\texttt{figures/hello.py} from synthetic data.

\begin{figure}[t]
  \centering
  \includegraphics[width=0.7\linewidth]{hello.pdf}
  \caption{A synthetic curve, regenerated on every build.}
  \label{fig:hello}
\end{figure}

\bibliographystyle{plainnat}
\bibliography{refs}
\end{document}
